
\chapter{طراحی Backend و ساختار Django}

\section{ساختار پیشنهادی}

\begin{latin}
\begin{lstlisting}
backend/
├── manage.py
├── config/
│   ├── settings.py
│   ├── urls.py
│   ├── asgi.py
│   └── wsgi.py
├── apps/
│   ├── accounts/
│   ├── content/
│   ├── creators/
│   ├── engagement/
│   ├── library/
│   ├── recommendations/
│   ├── moderation/
│   ├── messaging/
│   └── notifications/
└── tests/
\end{lstlisting}
\end{latin}

در عمل ممکن است Testهای هر App داخل همان App باشند. هدف این Diagram نشان دادن Domainهاست، نه اجبار به یک Layout خاص.

\section{config}

\filepath{config/} فقط Configuration پروژه است:
\begin{itemize}
  \item settings؛
  \item root URLs؛
  \item ASGI/WSGI؛
  \item environment-specific configuration later.
\end{itemize}

Business Model داخل config قرار نمی‌گیرد.

\section{accounts}

Modelهای احتمالی:
\begin{itemize}
  \item User با Avatar اختیاری و عمومی برای همهٔ حساب‌ها؛
  \item UserInterest؛
  \item UserPreference در صورت نیاز.
\end{itemize}

Serviceهای احتمالی:
\begin{itemize}
  \item update interests؛
  \item account data export/delete later.
\end{itemize}

\section{content}

هسته Domain:
\begin{itemize}
  \item Topic؛
  \item Source؛
  \item Idea؛
  \item IdeaTopic در صورت نیاز به Through Model.
\end{itemize}

قواعد:
\begin{itemize}
  \item فقط Published Idea در Feed عمومی؛
  \item Hidden content در Feed ظاهر نمی‌شود؛
  \item هر Idea باید Creator و Source مشخص داشته باشد؛
  \item Language باید مشخص باشد؛
  \item estimated reading time یا محاسبه شود یا قبل از publish ثبت شود.
\end{itemize}

\section{creators}

\begin{itemize}
  \item CreatorProfile؛
  \item CreatorFollow.
\end{itemize}

CreatorProfile با OneToOne به User وصل می‌شود.

\section{engagement}

\begin{itemize}
  \item Spark؛
  \item SavedIdea؛
  \item Comment.
\end{itemize}

Spark و Save باید Unique Constraint روی \eng{(user, idea)} داشته باشند.

\section{library}

در فاز اولیه Library می‌تواند صرفاً Query روی SavedIdea باشد.

\begin{decisionbox}
Table جدا برای LibraryEntry فقط زمانی ساخته می‌شود که Requirementی وجود داشته باشد که Save و Library Item را از هم متمایز کند.
\end{decisionbox}

بعداً:
\begin{itemize}
  \item reading stats؛
  \item eBooks؛
  \item membership entitlements.
\end{itemize}

\section{recommendations}

Recommendation Model اجباری نیست.

در ابتدا:
\begin{itemize}
  \item Query candidateها؛
  \item score در Service؛
  \item response.
\end{itemize}

اگر بعداً caching لازم شد، derived data طراحی می‌شود.

\section{moderation}

\begin{itemize}
  \item Report؛
  \item ModerationAction.
\end{itemize}

این App باید Permission boundary واضح داشته باشد.

\section{Business Logic کجا قرار بگیرد؟}

قاعده ساده:

\subsection{Model}

قواعد محلی و data invariants نزدیک خود Entity.

\subsection{Serializer}

Input/output transformation و validation API-level.

\subsection{View}

HTTP orchestration: request، permission، call business logic، response.

\subsection{Service}

Business operation چندمرحله‌ای یا قابل‌استفاده مجدد.

\subsection{Selector/Query helper}

Queryهای پیچیده read-side.

\section{Transaction}

Operationهایی که چند write وابسته دارند باید Transaction داشته باشند.

مثال future:
\begin{itemize}
  \item create moderation action + change target state؛
  \item membership activation + entitlement.
\end{itemize}

\section{Signals}

Django Signals جذاب‌اند اما باید محدود استفاده شوند.

\begin{warningbox}
Business Logic مهم را پشت Signalهای نامرئی پنهان نکن. اگر یک Action بخشی از Flow اصلی است، explicit service call معمولاً قابل‌فهم‌تر است.
\end{warningbox}

Signals برای side-effectهای محدود و کاملاً مستند مناسب‌اند.

\section{Settings}

ابتدا یک \filepath{settings.py} ساده کافی است. اگر بعداً Production/Development تفاوت زیاد شد می‌توان split کرد.

اصل:
\begin{itemize}
  \item Secret از env؛
  \item DEBUG در production false؛
  \item Database config از env؛
  \item Allowed hosts مشخص؛
  \item CORS/CSRF بر اساس deployment واقعی.
\end{itemize}

\section{Logging}

از \eng{print()} برای Debug موقت می‌توان استفاده کرد، اما Logging واقعی باید از Python logging/Django logging استفاده کند.

سطوح:
\begin{itemize}
  \item DEBUG؛
  \item INFO؛
  \item WARNING؛
  \item ERROR.
\end{itemize}

اطلاعات حساس نباید log شود.

\section{Admin strategy}

برای هر Model لازم نیست Admin بسیار پیچیده بسازیم.

ولی برای:
\begin{itemize}
  \item Idea؛
  \item Source؛
  \item Topic؛
  \item CreatorProfile؛
  \item Report
\end{itemize}
List display، search و filters مفید هستند.

\section{Seed data}

برای توسعه، Dataset کوچک و deterministic می‌سازیم:
\begin{itemize}
  \item چند User؛
  \item چند Creator؛
  \item 10-20 Topic؛
  \item چند Source از هر نوع؛
  \item حداقل 30-50 Idea برای Feed/Recommendation test.
\end{itemize}

این داده‌ها باید یا Fixture/Management Command قابل تکرار باشند، نه ورود دستی غیرقابل بازسازی.

\section{Management Command}

یک Command مثل:
\begin{latin}
\begin{lstlisting}
python manage.py seed_demo
\end{lstlisting}
\end{latin}
می‌تواند Demo data بسازد.

این بخش بعد از فهم basics اضافه می‌شود.

\section{Backend Definition of Done}

هر Backend feature:
\begin{itemize}
  \item Model/constraint درست؛
  \item Migration؛
  \item validation؛
  \item permission؛
  \item tests؛
  \item API docs؛
  \item admin/data setup در صورت لازم.
\end{itemize}
